\section{A uniform affine infrared tail from local cubic resolvents}
\label{sft-section}

This section bounds the large-physical-time contribution of the actual
soft scalar/vector pair. The proof retains the root block in every
trace and introduces no factor counting its \(K^4\) fine variables.
It uses the localized relative vertices proved above, not only a
global relative-form estimate. It does not determine the
intermediate-time logarithmic coefficient.

\subsection{Precise operator and source inputs}
\label{sft-input-section}

Write \(\mu=K^{-1}\). Fine scalar or positive-bond spaces use counting
inner products, with fixed color/component multiplicities. Assign a
bond to the block containing its initial site. Let \(P_x\) be the
orthogonal projection onto the final block \(x\), and put
\begin{equation}
 \mathcal L=I+K^2\Delta,
 \label{sft-reference}
\end{equation}
where \(\Delta\) is the ordinary componentwise flat lattice Laplacian.
Let \(X_0\) denote the flat physical operator, namely \(K^2\) times
the corresponding counting matrix.

For the vector spatial estimates, work on the chosen joint infinite
structural system or on explicit periodizations of its kernels.
Actual finite-regulator gaps and flat heat estimates are separate
facts. A finite momentum grid does not imply exact infinite moment
cancellation, and arbitrary finite winding kernels are not identified
with these periodizations. The following abstract proof is also
finite-volume uniform whenever its three inputs hold on the finite
system in question.

The first input is uniform positive ellipticity:
\begin{equation}
 X_0=X_0^*,\qquad X_0\ge c_0\mathcal L.
 \label{sft-ellipticity}
\end{equation}
The second is a relative physical-weight estimate. For every bounded
physical \(1\)-Lipschitz function \(f\), set \(W=e^{\theta f}\),
using the fine-site position divided by \(K\). On bonds use their
initial sites. For sufficiently small fixed \(\theta>0\), assume
\begin{equation}
 \left\|\mathcal L^{-1/2}
       (WX_0W^{-1}-X_0)\mathcal L^{-1/2}\right\|
       \le C\theta .
 \label{sft-weighted-input}
\end{equation}
Bounded truncations of distance functions suffice.

For the third input, let \(a_j\) be a complete physical \(C^1\) field
supported in a fixed smooth collar of block \(j\), and write
\begin{equation}
 a_j^*=\|a_j\|_{\mathrm{phys},1}
 =\|a_j\|_\infty+
 K\max_{\substack{\mathrm{nearest}\\\mathrm{same\ component}}}
       |a_j(e')-a_j(e)|.
 \label{sft-source-norm}
\end{equation}
Denote the first and mixed second variations of \(X\) at zero by
\(V_j\) and \(V_{jk}\). For every weight allowed above, require
\begin{align}
 \|\mathcal L^{-1/2}WV_jW^{-1}\mathcal L^{-1/2}\|
     &\le Ca_j^*,\notag\\
 \|\mathcal L^{-1/2}WV_{jk}W^{-1}\mathcal L^{-1/2}\|
     &\le Ca_j^*a_k^*.
 \label{sft-conjugated-vertices}
\end{align}
In addition, let \(f_j\) be physical distance to the fixed physical
cube or collar of block \(j\), and \(W_j=e^{\theta f_j}\), initially
with bounded truncations. The marked-source bounds are
\begin{align}
 \|\mathcal L^{-1/2}W_jV_jW_j\mathcal L^{-1/2}\|
     &\le Ca_j^*,\notag\\
 \|\mathcal L^{-1/2}W_jV_{jk}W_j\mathcal L^{-1/2}\|
     &\le Ca_j^*a_k^*,
 \label{sft-marked-vertices}
\end{align}
with the second bound also centered at \(k\).
The function \(f_j\) is Lipschitz on fine physical positions:
it is not a piecewise constant block-index distance. Its infimum
and supremum on any one block differ by a dimensional constant.

The two-sided weights in \eqref{sft-marked-vertices} attach a source
to both operator ends. An unlocalized bound on
\(X_0^{-1/2}V_jX_0^{-1/2}\) would not supply this input.
All source jets may be holomorphic and bilinear. Only the flat
operators \(X_0,\mathcal L\) are required to be positive.

\subsection{Why the actual soft pair supplies these inputs}
\label{sft-application-inputs}

For \(T_s=D_A^TD_A+\mu^2\Pi_A\), the soft scalar theorem gives
\eqref{sft-ellipticity}. Its weighted \(D,J\) expansion gives
\eqref{sft-weighted-input}. Its first two source derivatives give
\eqref{sft-conjugated-vertices}--\eqref{sft-marked-vertices}.
Here \(J\) is normalized and own-block local, so its derivatives
have no fine-fiber dimension loss. Scalar sources even need only
their supremum norm.

For the soft vector, the mean-zero Laplacian domination and gap,
preserved by the positive \(DJ\) update, imply
\eqref{sft-ellipticity} for \(X_0=K^2V_0\).
The relative-Hessian theorem gives, after physical rescaling,
first-two source jets of the form
\begin{equation}
 D_{\mathrm{phys}}^TZ+Z^TD_{\mathrm{phys}}+E,
 \label{sft-energy-vertex}
\end{equation}
with uniformly bounded, exponentially source-rooted factors \(Z,E\).
The second-order \(Z\) may have an additional beneficial factor
\(\mu\). Since
\(\|D_{\mathrm{phys}}\mathcal L^{-1/2}\|\le1\),
these factors give the relative vertex bounds. For example,
\[
 W_jD_{\mathrm{phys}}^TZW_j
 =D_{\mathrm{phys}}^T(W_jZW_j)
   +[W_j,D_{\mathrm{phys}}^T]W_j^{-1}(W_jZW_j).
\]
The commutator in the second term is \(O(\theta)\); source-rooted
decay bounds \(W_jZW_j\) at a smaller exponent. This proves the
two-sided normalization without commuting a fractional power of
\(\mathcal L\) through a weight.

The covariant \(D_AD_A^T\) term has the same elementary structure.
Each corrected angle average \(\bar C_A\), and each of its first
two physical-background derivatives, has counting size \(O(\mu)\).
Their pair therefore gives \(O(\mu^2)\), canceled by \(K^2\).
The coarse Hessian has fixed color fiber and uniform source bounds.
These terms supply zero-order vertices in
\eqref{sft-energy-vertex}. The soft-vector module separately
identifies this auxiliary family with the actual stationary affine
jets through order two, retaining the current and residual terms.

For completeness, the action contribution to
\eqref{sft-weighted-input} uses precisely the repaired ordered
Laurent factorization
\begin{equation}
 q_M=d_0^TH_0d_0.
 \label{sft-ordered-factorization}
\end{equation}
Both the constant and linear flat momentum jets vanish, and \(H_0\)
is symmetric with uniformly exponentially summable current kernels.
Separate component row sums alone would not justify this step;
no common componentwise difference-square identity is used.
Conjugation changes each outer gradient by a bounded multiplication
operator \(O(\theta\mu)\), and changes \(H_0\) by \(O(\theta\mu)\)
using its exponential first moment. Keep the latter change between
the two gradients. After multiplication by \(K^2\) and the two
\(\mathcal L^{-1/2}\) factors, the outer errors are
\(O(\theta+\theta^2)\) and the middle error is \(O(\theta\mu)\).
The longitudinal term obeys the same estimate. The bounded physical
block correction has fixed physical range, with a coarse Hessian
having uniform exponential moments, so its error is \(O(\theta)\).
This proves \eqref{sft-weighted-input} without a false uniform bound
on the absolute row norm of \(X_0\), which can grow as \(K^2\).

The uniform claims concern the explicit auxiliary source estimates.
Their identification with the actual affine coefficient is used at
each fixed \(K\); no uniform stationary-response bound is needed.

\subsection{A block Schatten bound without a rank loss}
\label{sft-schatten-section}

The elementary massive lattice heat estimate gives
\begin{equation}
 \operatorname{Tr}(P_x\mathcal L^{-3}P_x)\le C
 \label{sft-third-moment}
\end{equation}
uniformly in \(K\) and compatible volume. Indeed, integrate the
flat heat trace against \(t^2e^{-t}/2\). On a torus of current side
\(KN\), the heat trace per final block is at most
\[
 C\left\{\frac{K^4}{(1+K^2t)^2}+N^{-4}\right\}.
\]
Its indicated integral is bounded, including the finite zero mode.
The infinite lattice obeys the same bound without that mode.

Compression is not interchanged with taking powers. If \(u_i\) is
an orthonormal eigenbasis of \(P_x\mathcal L^{-1}P_x\) in the block,
scalar Jensen applied to the spectral measure of \(\mathcal L^{-1}\)
gives
\(\langle u_i,\mathcal L^{-1}u_i\rangle^3
 \le\langle u_i,\mathcal L^{-3}u_i\rangle\).
Therefore
\begin{equation}
 \|P_x\mathcal L^{-1/2}\|_{\mathfrak S_6}^6
 =\operatorname{Tr}[(P_x\mathcal L^{-1}P_x)^3]
 \le\operatorname{Tr}(P_x\mathcal L^{-3}P_x)\le C.
 \label{sft-s6}
\end{equation}
The analogous right-ended norm has the same bound. Schatten
Hölder with exponents \(6,\infty,6\) consequently yields
\begin{equation}
 \|P_x\mathcal L^{-1/2}A\mathcal L^{-1/2}P_y\|_{\mathfrak S_3}
 \le C\|A\|.
 \label{sft-s3}
\end{equation}
There is no factor proportional to the block rank \(K^4\).

\subsection{A common contour and rooted resolvent derivatives}
\label{sft-resolvent-section}

Choose \(0<\lambda_0<c_0/4\) and \(0<\varphi<\pi/2\).
Let \(\Gamma\) be the positively oriented Dunford contour on the
rays \(\lambda_0+r e^{\pm i\varphi}\), enclosing the positive
spectrum to their right. At finite regulator, first close it beyond
the largest eigenvalue. For \(R_z=(z-X_0)^{-1}\),
\begin{equation}
 \|\mathcal L^{1/2}R_z\mathcal L^{1/2}\|\le C_\Gamma,
 \qquad z\in\Gamma .
 \label{sft-contour-bound}
\end{equation}
Indeed, \(X_0\ge c_0\mathcal L\), while the spectral norm of
\(X_0^{1/2}R_zX_0^{1/2}\) is bounded on those rays.
By \eqref{sft-weighted-input} and a relative resolvent Neumann
expansion, choosing \(\theta\) sufficiently small gives
\begin{equation}
 WR_zW^{-1}=\mathcal L^{-1/2}A_{z,W}\mathcal L^{-1/2},
 \qquad \|A_{z,W}\|\le C ,
 \label{sft-twisted-resolvent}
\end{equation}
uniformly on the whole contour.

The first and mixed second derivatives are
\begin{align}
 \delta_jR_z&=R_zV_jR_z,\notag\\
 \delta_j\delta_kR_z
  &=R_zV_{jk}R_z+R_zV_jR_zV_kR_z+R_zV_kR_zV_jR_z.
 \label{sft-resolvent-derivatives}
\end{align}
Relative normalization cancels adjacent powers of \(\mathcal L\):
each product is \(\mathcal L^{-1/2}A\mathcal L^{-1/2}\)
with uniformly bounded \(A\).

Source localization can be checked without any fractional-weight
commutator. For a marked source \(j\),
\begin{equation}
 W_j(R_zV_jR_z)W_j
 =(W_jR_zW_j^{-1})(W_jV_jW_j)(W_j^{-1}R_zW_j).
 \label{sft-marked-first}
\end{equation}
For the ordered term with another source \(k\), put
\(R_+=W_jR_zW_j^{-1}\), \(R_-=W_j^{-1}R_zW_j\); then
\begin{equation}
 W_j(R_zV_jR_zV_kR_z)W_j
 =R_+(W_jV_jW_j)R_-(W_j^{-1}V_kW_j)R_-.
 \label{sft-marked-second}
\end{equation}
The marked vertex uses \eqref{sft-marked-vertices} and the other
vertex uses \eqref{sft-conjugated-vertices}. The primitive
\(V_{jk}\) has the marked bound at either source. Thus
\eqref{sft-s3} applies to every term.

Write \(\tau\) for the shortest connected graph length in the
final-block graph, allowing auxiliary vertices. The resulting
block bounds, for \(n=0,1,2\), are
\begin{equation}
 \|P_x\delta_{j_1}\cdots\delta_{j_n}R_zP_y\|_{\mathfrak S_3}
 \le C_n\prod_{\ell=1}^n a_{j_\ell}^*
       e^{-c\tau(x,y,j_1,\ldots,j_n)}.
 \label{sft-rooted-resolvent}
\end{equation}
For \(n=0\), use a one-sided distance weight from \(y\).
For \(n=1\), the marked weight gives
\(e^{-\theta[d(x,j)+d(y,j)]}\). For \(n=2\), take the
geometric mean of the separate \(j\)- and \(k\)-centered estimates.
Their four edges connect all marked vertices and dominate a
connecting graph length. The infimum/supremum comparison of
physical distances within each block changes only the constant.
Fixed collars have the same effect. Periodic physical distances
and bounded truncations give the identical conclusion whenever the
inputs hold; the uniform bounds permit removal of the truncations.

\subsection{Three resolvents preserve the local trace}
\label{sft-heat-section}

Functional calculus gives, for \(u>0\),
\begin{equation}
 e^{-uX}=\frac{2}{u^2}\frac1{2\pi i}
       \int_\Gamma e^{-uz}(z-X)^{-3}\,dz .
 \label{sft-cubic-calculus}
\end{equation}
The scalar residue of the third-order pole is
\(u^2e^{-u\lambda}/2\), fixing both coefficient and sign.

Differentiate \(R_z^3\) through order two. Each term is a product of
three differentiated resolvents, with derivative orders
\(n_1+n_2+n_3\le2\). Inserting the block resolution of the identity
between them gives
\begin{align}
 &\left|\operatorname{Tr}
  (P_0R_z^{(n_1)}P_iR_z^{(n_2)}P_jR_z^{(n_3)}P_0)\right|
 \notag\\
 &\qquad\le
  \|P_0R_z^{(n_1)}P_i\|_{\mathfrak S_3}
  \|P_iR_z^{(n_2)}P_j\|_{\mathfrak S_3}
  \|P_jR_z^{(n_3)}P_0\|_{\mathfrak S_3}.
 \label{sft-local-holder}
\end{align}
Here superscripts denote the indicated source derivatives, not
ordinary powers. The root block stays in place. Nothing is moved
cyclically through \(P_0\).

Equation \eqref{sft-rooted-resolvent} and exponential convolution,
with a smaller exponent, imply
\begin{equation}
 \left|\delta_{j_1}\cdots\delta_{j_n}
       \operatorname{Tr}(P_0R_z^3P_0)\right|
 \le C_n\prod_{\ell=1}^n a_{j_\ell}^*
        e^{-c'\tau(0,j_1,\ldots,j_n)},\qquad n\le2.
 \label{sft-cubic-root}
\end{equation}
The two intermediate block sums are absolutely convergent in trace
norm, uniformly in \(K,z\) and the applicable volume. This constructs
the infinite local trace without an infinite extensive trace.

On the contour,
\(\operatorname{Re}z=\lambda_0+r\cos\varphi\), whence
\[
 \int_\Gamma |e^{-uz}|\,|dz|
       \le C u^{-1}e^{-\lambda_0u}.
\]
Combining this with
\eqref{sft-cubic-calculus}--\eqref{sft-cubic-root} proves
\begin{equation}
 \left|\delta_{j_1}\cdots\delta_{j_n}
       \operatorname{Tr}(P_0e^{-uX}P_0)\right|
 \le C_nu^{-3}e^{-\lambda_0u}
       \prod_{\ell=1}^n a_{j_\ell}^*
       e^{-c'\tau(0,j_1,\ldots,j_n)},\quad u\ge1,\ n\le2.
 \label{sft-rooted-heat}
\end{equation}
The displayed absolute majorants justify the differentiations,
contour integral, block sums and local limits. At fixed \(K\), begin
with bounded compact sources in the actual local germ, then extend
their coefficients using these bounds. No small-time asymptotic is
claimed by the harmless factor \(u^{-3}\).

\subsection{The actual affine infrared contribution}
\label{sft-affine-section}

Choose a smooth physical partition of unity on fixed block collars,
with uniformly bounded first derivatives and finite overlap.
Decompose the affine physical field
\begin{equation}
 A_\alpha(x)=\frac{x_\beta}{K}H,\qquad \alpha\ne\beta,
 \label{sft-affine}
\end{equation}
into complete supported fields \(a_j\).
Their norms satisfy \(a_j^*\le C(1+|j|)\). The second source bound
\eqref{sft-rooted-heat} is summable with these polynomial weights.
Consequently
\begin{equation}
 \left|[t^2]\operatorname{Tr}_0 e^{-uX(tA)}\right|
       \le Ce^{-cu},\qquad u\ge1,
 \label{sft-affine-heat}
\end{equation}
uniformly in \(K\) and old depth on the stated structural systems.
This is evaluation of exponentially localized Taylor kernels,
not insertion of a nonzero unbounded field into a fixed analytic ball.

The actual first two stationary affine jets agree with the explicit
representatives just estimated, up to a fine gauge. Their diagonal
color traces are individually gauge invariant. Thus the fixed-\(K\)
polynomial-jet identification transfers
\eqref{sft-affine-heat} to the actual coefficient before taking the
uniform bound. It does not require a uniform-\(K\) norm for the
full stationary response.

For \(X_s=K^2T_s\) and \(X_v=K^2V\), we obtain
\begin{theorem}[Uniform actual infrared tail]
\label{sft-main}
Under the established soft-factor and localized relative-vertex
inputs, the actual affine coefficient satisfies
\begin{equation}
 \int_1^\infty
  \left|[t^2]\left\{
   \operatorname{Tr}_{0,E}e^{-uX_v(t)}
     -2\operatorname{Tr}_{0,S}e^{-uX_s(t)}
                 \right\}\right|\frac{du}{u}\le C .
 \label{sft-tail}
\end{equation}
The constant is uniform in \(K\), old depth and the chosen joint
infinite structural systems. Each factor separately obeys the same
type of bound.
\end{theorem}

The fine link angle corresponding to \eqref{sft-affine} is
\(t x_\beta H/K^2\). The change of variables \(s=K^2u\)
identifies \eqref{sft-tail} with the current-time tail \(s\ge K^2\)
of the actual soft-pair Maxwell heat formula.
Constant determinant normalizations have zero affine coefficient.

No finite torus affine-flux family is asserted. Finite source bounds
apply to explicit periodizations or to finite models for which the
three inputs were separately verified. The actual polynomial
coefficient is obtained in the chosen infinite structural system by
the fixed-\(K\) identification; unrelated actual finite winding
kernels are not silently substituted.

Together with the exact soft factorization and the bounded coarse
Maxwell remainders, \eqref{sft-tail} removes the large-physical-time
part from the logarithmic calculation. The remaining comparison is
the short-current-time contribution and the intermediate range
\(1\ll s\ll K^2\). Neither the value nor the sign of the logarithmic
coefficient follows from this tail estimate.
